% ===== PRÉAMBULE CANONIQUE — à coller VERBATIM en tête de CHAQUE .tex =====
\documentclass[11pt,a4paper]{article}
\usepackage[utf8]{inputenc}\usepackage[T1]{fontenc}\usepackage{lmodern}
\usepackage[french]{babel}
\usepackage{amsmath,amssymb,mathtools}\usepackage{icomma}
\usepackage[version=4]{mhchem}          % \ce{...} : équations & formules
\usepackage{siunitx}\sisetup{locale=FR} % \qty{}{}, \si{}, \ang{}
\usepackage{chemfig}                    % \chemfig{...} : structures développées
\usepackage{xcolor}\usepackage{colortbl}
\usepackage{tikz}\usepackage{pgfplots}\pgfplotsset{compat=1.18}
\usepackage[most]{tcolorbox}\usepackage{enumitem}\usepackage{array,booktabs}
\usepackage[a4paper,margin=2.2cm]{geometry}
\usetikzlibrary{arrows.meta,calc,angles,quotes,positioning,patterns,backgrounds,shapes.geometric}
\definecolor{primary}{RGB}{222,49,99}\definecolor{primarydark}{RGB}{150,22,55}
\definecolor{defcol}{RGB}{0,114,178}\definecolor{methcol}{RGB}{52,92,148}
\definecolor{excol}{RGB}{0,135,90}\definecolor{attncol}{RGB}{200,40,55}
\tcbset{boxbase/.style={enhanced,breakable,boxrule=0.9pt,arc=2.5pt,
  left=3mm,right=3mm,top=2.3mm,bottom=2.3mm,fonttitle=\bfseries,coltitle=white}}
% --- boîtes TUTORIEL (noms EXACTS, ne pas renommer) ---
\newtcolorbox{cle}[1][]{boxbase,colback=defcol!6,colframe=defcol,
  title={\textsc{À retenir}\ifx\relax#1\relax\relax\else\textmd{ : \itshape #1}\fi}}
\newtcolorbox{methode}[1][]{boxbase,colback=methcol!7,colframe=methcol,sharp corners,
  title={\textsc{Méthode}\ifx\relax#1\relax\relax\else\textmd{ --- \itshape #1}\fi}}
\newtcolorbox{exemplebox}[1][]{boxbase,colback=excol!5,colframe=excol,
  title={\textsc{Exemple}\ifx\relax#1\relax\relax\else\textmd{ : \itshape #1}\fi}}
\newtcolorbox{piege}[1][]{boxbase,colback=attncol!6,colframe=attncol,
  title={\textsc{Piège}\ifx\relax#1\relax\relax\else\textmd{ : \itshape #1}\fi}}
\newtcolorbox{remarque}[1][]{boxbase,colback=black!4,colframe=black!45,
  title={\textsc{Remarque}\ifx\relax#1\relax\relax\else\textmd{ : \itshape #1}\fi}}
% --- boîtes EXERCICE / CORRIGÉ (noms EXACTS, auto-comptées) ---
\newtcolorbox[auto counter]{exo}[1][]{boxbase,colback=primary!4,colframe=primary,
  title={\textsc{Exercice}~\thetcbcounter\ifx\relax#1\relax\relax\else\textmd{ --- \itshape #1}\fi}}
\newtcolorbox[auto counter]{corrige}[1][]{boxbase,colback=excol!4,colframe=excol,
  title={\textsc{Corrigé}~\thetcbcounter\ifx\relax#1\relax\relax\else\textmd{ --- \itshape #1}\fi}}
\newcommand{\R}{\mathbb{R}}\newcommand{\N}{\mathbb{N}}\newcommand{\Z}{\mathbb{Z}}\newcommand{\ds}{\displaystyle}
% =========================================================================
\begin{document}
\begin{exo}[écrire la constante d'équilibre]
Écris l'expression de la constante d'équilibre $K$ (en concentrations).
\begin{enumerate}[label=\alph*)]
  \item $\ce{N2}(g) + 3\,\ce{H2}(g) \rightleftharpoons 2\,\ce{NH3}(g)$
  \item $\ce{H2}(g) + \ce{I2}(g) \rightleftharpoons 2\,\ce{HI}(g)$
  \item $\ce{PCl5}(g) \rightleftharpoons \ce{PCl3}(g) + \ce{Cl2}(g)$
\end{enumerate}
\end{exo}

\begin{exo}[quelles espèces figurent dans $K$ ?]
Écris $K$ en \textbf{excluant} les solides et liquides purs.
\begin{enumerate}[label=\alph*)]
  \item $\ce{CaCO3}(s) \rightleftharpoons \ce{CaO}(s) + \ce{CO2}(g)$
  \item $2\,\ce{A}(aq) + \ce{B}(s) \rightleftharpoons 2\,\ce{C}(g) + \ce{D}(aq)$
\end{enumerate}
\end{exo}

\begin{exo}[interpréter la valeur de $K$]
Vers quoi l'équilibre est-il déplacé ?
\begin{enumerate}[label=\alph*)]
  \item $K = 10^{5}$
  \item $K = 10^{-4}$
  \item $K \approx 1$
\end{enumerate}
\end{exo}

\begin{exo}[vrai ou faux : l'état d'équilibre]
Indique \textbf{vrai} ou \textbf{faux}.
\begin{enumerate}[label=\alph*)]
  \item « À l'équilibre, les concentrations restent constantes dans le temps. »
  \item « À l'équilibre, les réactions directe et inverse s'arrêtent totalement. »
  \item « À l'équilibre, réactifs et produits sont nécessairement en quantités égales. »
  \item « La valeur de $K$ dépend de la température. »
\end{enumerate}
\end{exo}

\begin{exo}[calculer $K$ en une ligne]
Le système $\ce{PCl5}(g) \rightleftharpoons \ce{PCl3}(g) + \ce{Cl2}(g)$ est à l'équilibre avec $[\ce{PCl5}]=\qty{0.40}{\mol\per\litre}$, $[\ce{PCl3}]=\qty{1.40}{\mol\per\litre}$ et $[\ce{Cl2}]=\qty{0.50}{\mol\per\litre}$.
\begin{enumerate}[label=\alph*)]
  \item Écris $K$.
  \item Calcule sa valeur et conclus (vers réactifs ou produits ?).
\end{enumerate}
\end{exo}

\end{document}
